% Slide "What a data stream is": the memory one step of the inner loop touches (AB pattern, the
% generated StreamCollide kernel). Each of the 19 values of a cell has its own array (waLBerla's
% fzyx layout); pull streaming reads value i at the neighbour x - c_i (still consecutive addresses
% along x), the AB pattern writes a second set of 19 arrays, and the kernel writes rho and u
% (4 arrays): 19 + 19 + 4 = 42 sequences of consecutive addresses, all advancing together.
% Evidence: the generated StreamCollide of SweepgenLdcCPU_explicit_cse_pf (GCC f32, board) has
% exactly 19 prefetch_read and 23 prefetch_write calls in its loop, one per array it walks through.
% presentation5: animated. Overlays 1-3: one array only, the loop moves on by one cache line per overlay
% (one 64-byte line = 16 cells in f32 = 2 RVV loads), used lines turn grey, the hardware prefetcher has the
% next lines on their way; 4: all 42 arrays, advancing together.
\begin{tikzpicture}[x=1cm,y=1cm]
  \tikzset{cl/.style={draw=psgray!70,fill=white,line width=0.3pt},
           clnow/.style={draw=psteal!80!black,fill=psteal!65,line width=0.3pt},
           clnext/.style={draw=psteal!60,fill=psteal!18,line width=0.3pt},
           cldone/.style={draw=psgray!70,fill=psgray!28,line width=0.3pt},
           arrlbl/.style={font=\scriptsize,anchor=east,inner sep=1pt}}
  %   \lines{y}{current line}{extra style}: one array as 10 cache lines; lines before the current one are done,
  %   the two after it are on their way
  \newcommand{\lines}[3]{%
    \foreach \k in {0,...,9} {
      \ifnum\k<#2 \def\st{cldone}\else\ifnum\k=#2 \def\st{clnow}\else
        \ifnum\k<\numexpr#2+3\relax \def\st{clnext}\else \def\st{cl}\fi\fi\fi
      \draw[\st,#3] (\k*0.4,#1) rectangle ++(0.4,0.22);
    }}
  %   \strip{y}{label}{visible from}: the array moves with the loop on overlays 1-3, then stays at line 5
  \newcommand{\strip}[3]{%
    \only<1>{\lines{#1}{3}{visible on=<#3->}}\only<2>{\lines{#1}{4}{visible on=<#3->}}%
    \only<3->{\lines{#1}{5}{visible on=<#3->}}%
    \node[arrlbl,visible on=<#3->] at (-0.05,#1+0.11) {#2};}
  \newcommand{\vdotsat}[1]{\node[font=\scriptsize,text=psgray,visible on=<4->] at (2.2,#1) {$\vdots$};}
  \newcommand{\grp}[4]{% y bottom, y top, count, what
    \draw[decorate,decoration={brace,amplitude=3pt},psgray!90!black,visible on=<4->] (4.15,#2) -- (4.15,#1)
      node[midway,right=3pt,font=\scriptsize,align=left] {\textbf{#3} #4};}

  \strip{3.5}{$f_0$}{1} \strip{3.23}{$f_1$}{4} \vdotsat{2.98} \strip{2.62}{$f_{18}$}{4}
  \grp{2.62}{3.72}{19}{read}
  \strip{2.12}{$f'_0$}{4} \strip{1.85}{$f'_1$}{4} \vdotsat{1.6} \strip{1.24}{$f'_{18}$}{4}
  \grp{1.24}{2.34}{19}{written}
  \strip{0.74}{$\rho$}{4} \strip{0.47}{$u_x$}{4} \vdotsat{0.22} \strip{-0.14}{$u_z$}{4}
  \grp{-0.14}{0.96}{4}{outputs}

  %   1-3: where the loop is, in the one array shown
  \foreach \o/\k in {1/3,2/4,3/5}
    \node[font=\tiny,text=psteal!50!black,anchor=north,inner sep=1pt,visible on=<\o>] at (\k*0.4+0.2,3.47) {$\uparrow$ loop};
  %   in the space the other arrays take from overlay 4 on
  \node[font=\scriptsize,anchor=north west,align=left,text width=4.2cm,visible on=<1-3>] at (-0.05,3.15)
    {\hyphenpenalty=10000\exhyphenpenalty=10000 Every 16 cells (in f32) the loop reaches the next cache line of 64 bytes. The hardware prefetcher
     has recognised the pattern and is \mbox{already} loading the next lines.};

  \draw[arr] (0,-0.42) -- (4.0,-0.42) node[midway,below,lbl] {the loop goes along $x$: addresses grow};
  \node[draw=psteal!80!black,fill=psteal!8,rounded corners=2pt,font=\small\bfseries,anchor=west,visible on=<4->]
    at (4.3,-0.45) {= 42 streams};
  %   legend
  \draw[clnow] (0,3.95) rectangle ++(0.3,0.18);
  \node[lbl,anchor=west] at (0.35,4.04) {in use};
  \draw[clnext] (1.35,3.95) rectangle ++(0.3,0.18);
  \node[lbl,anchor=west] at (1.7,4.04) {loading early};
  \draw[cldone] (3.65,3.95) rectangle ++(0.3,0.18);
  \node[lbl,anchor=west] at (4.0,4.04) {already used};
\end{tikzpicture}
